\documentclass[11pt]{article}
\usepackage[margin=1in]{geometry}
\usepackage{booktabs}
\usepackage{amsmath}
\usepackage{hyperref}
\usepackage{natbib}
\usepackage{xcolor}
\title{Understand the User, Don't Repeat:\\
Linguistic Cues of Unmet Needs and Explicit Repair in Task-Oriented Dialogue}
\author{Ying Chen}
\date{}

\begin{document}
\maketitle

\begin{abstract}
When a dialogue system fails, users seldom say so in emotional words. They repeat the request, rephrase it,
write ``what?'', or get shorter. We study these cues in EmoWOZ, a corpus with emotion labels on every user turn.
Human operators change their replies after a dissatisfied turn (acknowledgment rises from 3.9\% to 17.8\%);
a deployed machine system does not (3.7\% vs.\ 3.5\%). Toward the machine, 46.8\% of dissatisfied turns repeat
an earlier request, and only 7.7\% contain a complaint or negation word, so word-level sentiment misses most
of them. We then ask whether a language model can be made to answer the unmet need rather than the last
message. Asking the model to ``read the user first'' helps a strong model but hurts an 8B model, which
misreads the need and apologizes instead of delivering. We propose REPAIR, which reads the dialogue
explicitly --- repair cues, a ledger of requested, delivered and pending items, and database rows for the
pending item --- and hands this reading to the generator. On 100 unseen dialogue cases, REPAIR raises how
well an 8B model resolves the user's need from 0.45 to 1.22 on a 0--2 scale (0 = unresolved, 1 = partly,
2 = delivered), cuts the share of replies after which the user would likely give up from 49\% (prompt-only
reading) to 19\%, and brings the 8B model to the level of a much stronger model given a ``be concrete''
instruction.
\end{abstract}

\section{Introduction}
A user who has asked twice for a guesthouse in the south does not need a third question about the price
range, and does not need an apology. The user needs the guesthouse. Dialogue systems, including systems built
on large language models, are usually asked to respond to the last message. When an earlier request is still
open, the last message is often a repetition (``hotel'', ``hotel options'', ``four star hotels''), and a system
that answers it in isolation repeats its own failure.

Conversation analysis calls these moves \emph{other-initiated repair}: the listener signals that something
was not understood or not delivered and invites the speaker to fix it \citep{schegloff1977repair}. In the
terms of \citet{clark1991grounding}, the request has not been grounded. Our view is that unmet needs are
visible mostly in this linguistic behaviour --- repetition, rephrasing, repair tokens, shortening --- and only
partly in emotional vocabulary. Earlier work by the author on offensive language took the same position:
lexical and syntactic structure, together with a user's writing style, carried most of the signal
\citep{chen2012offensive}.

This paper makes three contributions.
(1) A corpus study on EmoWOZ \citep{feng2022emowoz} and DBDC3 \citep{higashinaka2016dbdc} showing that a
machine system does not react to dissatisfaction, that users express dissatisfaction in task dialogue mainly
by climbing a repair ladder (dissatisfaction rises from 16.3\% at stage~0 to 79.3\% at stage~3+, capturing
81.7\% of negative turns vs.\ 7.7\% for complaint words), and that repetition is task-specific.
(2) A negative result: telling a model to infer the user's state from the dialogue helps a strong model and
hurts a small one.
(3) REPAIR, an explicit reading of repair cues and pending needs, grounded in the task database, which makes
a small model resolve the user's need far more often on unseen dialogues and outperforms ledger-only grounding
on high-escalation turns. We release code, prompts and all model outputs.

\section{Related Work}
\paragraph{Emotion in task-oriented dialogue.} EmoWOZ labels every user turn of MultiWOZ
\citep{budzianowski2018multiwoz} and of 1,000 human--machine dialogues (DialMAGE) with one of seven emotions
\citep{feng2022emowoz}. Recent work detects frustration in deployed systems
\citep{frustration2024stupidrobot} and studies emotion in customer-service dialogue, finding that apology
and empathy often leave anger in place \citep{emowozcs2025}. These works detect the state; we study how it
is expressed and how to act on it.
\paragraph{Repair and grounding.} Our cue set follows the repair literature \citep{schegloff1977repair} and
grounding theory \citep{clark1991grounding}.
\paragraph{Dialogue breakdown detection.} The DBDC challenge series \citep{higashinaka2016dbdc} annotates
system turns as breakdowns in chat-oriented dialogue. We use DBDC3 as a contrast corpus: without task goals,
user repetition barely rises after breakdowns.
\paragraph{Dialogue state tracking.} Classical dialogue state trackers infer slot--value belief states across
turns \citep{budzianowski2018multiwoz}. REPAIR is a lightweight, zero-shot alternative focused on \emph{repair
and grounding state}: rather than training a neural tracker, it pairs surface repair cues with a ledger of
requested versus delivered items so any off-the-shelf generator can close pending needs without schema-specific
fine-tuning.
\paragraph{LLM judges.} LLM judges have known biases \citep{zheng2023judging}; we report one more in task
dialogue: they reward confident, invented bookings.

\section{Corpus Study}
\label{sec:corpus}
EmoWOZ has two parts: MultiWOZ, where a human ``wizard'' plays the system (10,438 dialogues), and DialMAGE,
where a machine system talks to people (995 dialogues). We call a user turn \emph{negative} if it is labelled
dissatisfied or abusive. Confidence intervals are 95\% bootstrap intervals over dialogues (2,000 resamples).

\paragraph{The machine does not react.} We match apology and acknowledgment words (\emph{sorry, apologize,
my mistake, I understand}, \ldots) in the system's next reply. Human wizards use them after 17.8\% of negative
turns and 3.9\% of neutral turns. The machine: 3.5\% and 3.7\%. Including \emph{unfortunately / I'm afraid}:
wizards 20.3\% [17.7, 22.7] vs.\ 5.4\%; machine 5.0\% [4.2, 5.8] vs.\ 5.6\% [5.0, 6.3].

\paragraph{Dissatisfaction is expressed by repeating.} Table~\ref{tab:cues} shows, for each negative turn,
whether it contains a complaint or negation word, and whether it repeats an earlier request of the same user
(token Jaccard $\ge 0.5$).

\begin{table}[h]
\centering
\begin{tabular}{lrrrr}
\toprule
 & $n$ & complaint word & repeats own request & neither \\
\midrule
MultiWOZ, negative turns & 958 & 53.0\% & 3.0\% & 44.6\% \\
DialMAGE, negative turns & 4,264 & 7.7\% & 46.8\% & 47.6\% \\
DialMAGE, neutral turns & 7,230 & 3.0\% & 13.2\% & 84.2\% \\
\bottomrule
\end{tabular}
\caption{How dissatisfaction is expressed. With a human wizard users complain in words; with a machine
they repeat themselves.}
\label{tab:cues}
\end{table}

When the machine repeats its own previous reply (Jaccard $\ge 0.8$) after a negative turn, the next user
turn is still negative 88.3\% [80.0, 95.7] of the time, against 54.6\% [52.7, 56.6] otherwise ($n=71$
repeats).

\paragraph{Is repetition specific to task dialogue, and does the finding hold on other task corpora?} In the
DBDC3 chat-oriented corpus \citep{higashinaka2016dbdc} (415 dialogues, no task goals), users repeat after a
system breakdown only 8.4\% [6.5, 10.4] of the time, barely above the 6.8\% [5.9, 7.8] baseline ($+1.5$ pp).
In DialMAGE the gap is $+33.6$ pp: users repeat when they have an unmet \emph{task} need, not simply when the
system is incoherent. To verify that our cue findings are not specific to MultiWOZ domains, we apply the same
zero-shot repair reader (\texttt{analyze\_uss.py}) to two external goal-oriented benchmarks annotated with
turn-level user satisfaction \citep{sun2021uss}: Schema-Guided Dialogue \citep[SGD;][]{rastogi2020sgd} (1,000
dialogues across 20 domains such as banks, flights, media and doctors; 9,833 user turns, 2,194 dissatisfied)
and ReDial movie recommendation (999 dialogues; 7,619 user turns, 1,793 dissatisfied). Explicit emotion words
(\emph{frustrated, annoyed, useless, angry}, \ldots) appear in only \textbf{0.3\%} [0.1, 0.5] of dissatisfied
turns in SGD and \textbf{2.7\%} [1.9, 3.6] in ReDial, whereas REPAIR's surface repair cues (stage $\ge 1$)
capture \textbf{63.6\%} [61.4, 65.8] and \textbf{51.2\%} [48.9, 53.8], doubling the dissatisfaction rate over
stage~0 in SGD (30.5\% vs.\ 15.2\%) and raising it by $+12.3$ pp in ReDial (31.1\% vs.\ 18.8\%).

These findings motivate a method that detects repetition-based repair cues and tracks pending needs, rather
than relying on sentiment words or the model's own reflection.

\section{REPAIR}
\label{sec:method}
REPAIR (Reading Escalation And Pending Items for Response) builds a short, checkable reading of the
dialogue at the current user turn. It is deterministic and uses no model calls.
\begin{enumerate}
\item \textbf{Repair reader.} Cues in the last user turn: repetition (Jaccard $\ge 0.5$ with an earlier user
turn) or rephrasing ($0.3$--$0.5$); repair tokens (\emph{what?}, \emph{??}, \emph{I said}, \emph{already});
negation (\emph{no, not, again, still}); talk about the conversation (\emph{can you help}, \emph{understand});
a turn under half the user's median length; give-up words. These map to an escalation stage 0--4.
\item \textbf{Grounding ledger.} For each domain (restaurant, hotel, attraction, train, taxi, \ldots), the
constraints the user stated, whether a matching entity was delivered by the system afterwards, and which
required details are missing (train: origin, destination, day; taxi: origin, destination, time). Turns
without domain words refine the current domain; a bare answer fills the slot the system asked for; negated
needs (``does not need wifi'') remove a constraint. Only the current topic is shown.
\item \textbf{Database grounding.} Rows from the MultiWOZ database that match the pending request; if none,
constraints are relaxed one at a time (type and amenities first) and the relaxation is stated.
\item \textbf{Policy.} By stage: deliver the pending item first; ask only for missing required details; book
when asked; at stage 4, say what can and cannot be done and offer a human.
\end{enumerate}
The reading is appended to the dialogue, together with the system's last turns marked ``do not repeat''.
Any generator can use it. Among REPAIR's four components, the repair reader and the escalation policy are
completely domain-agnostic (as validated zero-shot on SGD's 20 domains above); for the grounding ledger and
database lookup, all entity names, categorical slot values (e.g.\ food types, attraction types, stations) and
searchable columns are induced directly from the target JSON database tables at startup rather than
hand-listed per venue, so adapting to a new schema requires only supplying its database table and a short map
of domain trigger words.

\paragraph{The repair ladder tracks emotion across the corpus.} Applying REPAIR's stage reader and grounding
ledger to all 12,093 annotated user turns in DialMAGE (\texttt{analyze\allowbreak\_ladder.py}), user dissatisfaction
rises monotonically up the repair ladder (Table~\ref{tab:ladder}): from 16.3\% at stage~0 (no repair cues) to
34.3\% at stage~1 (rephrasing, negation or shortening), 60.3\% at stage~2 (repeating an earlier request or
other-initiated repair), and 79.3\% at stage~3+ (repeated repair or meta-complaint), capturing 81.7\% of all
negative turns at stage~$\ge 1$ (vs.\ 7.7\% for complaint words). Next-turn dissatisfaction likewise climbs
from 32.7\% after stage~0 to 57.1\% after stage~3+. Independently, turns where the grounding ledger shows an
undelivered request from an earlier turn are negative 40.5\% [38.9, 42.1] of the time ($n=7{,}919$) versus
25.3\% [23.3, 27.3] otherwise ($n=4{,}174$).

\begin{table}[h]
\centering
\begin{tabular}{lrrr}
\toprule
Repair ladder stage & $n$ & current negative & next turn negative \\
\midrule
Stage 0 (no repair cues) & 4,767 & 16.3\% [15.0, 17.7] & 32.7\% [31.1, 34.2] \\
Stage 1 (rephrase, negation, shorter) & 4,242 & 34.3\% [32.5, 36.1] & 37.9\% [36.0, 39.7] \\
Stage 2 (repeat $1\times$, \emph{what?}, meta) & 2,177 & 60.3\% [58.0, 62.5] & 45.3\% [42.8, 47.8] \\
Stage 3+ (repeat $\ge 2\times$ or + OIR/meta) & 907 & \textbf{79.3\%} [75.7, 82.7] & \textbf{57.1\%} [52.8, 61.4] \\
\bottomrule
\end{tabular}
\caption{User dissatisfaction by repair escalation stage across all 12,093 DialMAGE user turns (95\% bootstrap
CIs over dialogues). The emotion label follows the repair ladder rather than emotional vocabulary.}
\label{tab:ladder}
\end{table}

\section{Experiments}
\paragraph{Cases.} From DialMAGE we take every dialogue at the first user turn that is negative and repeats
an earlier request: the need is still unmet. We developed prompts and the reader on 60 cases of the test
split, and on 81 cases of the dev split for a second round; the results we report are on 100 cases from the
train split, which were not looked at during development.
\paragraph{Conditions.} A: reply to the last message. B: be concrete, answer rather than ask back.
C: prompt-only reading --- find what is unresolved from the same linguistic cues and how the user feels,
then deliver it (for the 8B model, two calls: read, then reply). D: REPAIR. E: REPAIR without stage, cues and
policy (ledger and database only).
\paragraph{Generators.} Claude 3.5 Sonnet \citep{anthropic2024claude} via the
\texttt{claude} CLI, and Llama 3 8B \citep{grattafiori2024llama3} (dolphin-llama3:8b
via Ollama). A--C have no database; D and E have the rows REPAIR selects.
\paragraph{Evaluation.} Two judge models from different families (Claude 3 Opus and Gemini 2.5
\citep{geminiteam2024gemini}) independently see
the dialogue and the five replies in random order, and score each reply: \emph{resolves} (0 does not, 1
partly, 2 delivers the need), \emph{give\_up} (a user in this state would likely give up), \emph{repeats},
and the best reply. In a second pass each judge also sees the database rows for the current request and is
told that a contradicting fact cannot score 2. Across all 500 judged replies per run, the two judge families
agree within 1 point on \emph{resolves} 96.4\%--100.0\% of the time; fact-checking raises inter-judge
correlation from $r = 0.73$ to $r = 0.90$ on the 8B run (exact agreement 59.4\% to 85.2\%) and from $r = 0.66$
to $r = 0.81$ on the Sonnet run. An objective check without any model asks whether the reply
names a real database entity that satisfies the pending constraints. Differences are paired bootstrap 95\%
intervals over cases.

\subsection{Results}
\begin{table}[h]
\centering
\begin{tabular}{lccccc}
\toprule
 & \multicolumn{2}{c}{Claude Opus judge} & \multicolumn{2}{c}{Gemini judge} & \\
\cmidrule(lr){2-3} \cmidrule(lr){4-5}
Llama 3 8B & resolves & give up & resolves & give up & real \& match \\
\midrule
A last message & 0.45 & 29\% & 0.98 & 19\% & 23.0\% \\
B be concrete & 0.45 & 28\% & 0.98 & 18\% & 21.8\% \\
C prompt-only reading & 0.33 & 49\% & 0.73 & 35\% & 18.4\% \\
D REPAIR & \textbf{1.22} & \textbf{19\%} & \textbf{1.51} & \textbf{14\%} & 60.9\% \\
E ledger + database & 1.16 & 23\% & 1.39 & 22\% & 62.1\% \\
\midrule
A, fact-checked judge & 0.20 & 54\% & 0.26 & 41\% & \\
B, fact-checked judge & 0.17 & 56\% & 0.26 & 39\% & \\
C, fact-checked judge & 0.20 & 69\% & 0.25 & 51\% & \\
D, fact-checked judge & \textbf{1.41} & \textbf{14\%} & \textbf{1.49} & \textbf{12\%} & \\
E, fact-checked judge & 1.31 & 17\% & 1.39 & 16\% & \\
\bottomrule
\end{tabular}
\caption{100 unseen DialMAGE dialogues, Llama 3 8B generator, judged independently by Claude Opus and Gemini
(both different model families from the generator). ``Real \& match'' is over the 87 cases with a database
request (the original DialMAGE system achieved 10.3\%, after which 53.6\% of users stayed negative).}
\label{tab:llama}
\end{table}

\paragraph{Small model.} Under both judges, REPAIR resolves the need far more often than every baseline
(Table~\ref{tab:llama}): with Claude Opus, D$-$B $+0.77$ [0.55, 0.98] and D$-$C $+0.89$ [0.70, 1.08]; with
Gemini, D$-$B $+0.53$ [0.33, 0.72] and D$-$C $+0.78$ [0.60, 0.96]. Under both judges, prompt-only reading
(C) significantly hurts the 8B model compared to B (Opus $-0.12$, Gemini $-0.25$ [$-0.42$, $-0.08$], give-ups
$+17$ pts): inspecting 18 development cases where C lost, the model's own reading named needs from earlier in
the dialogue or invented them (``Indian restaurant'' for a guesthouse request), and the reply apologized and
promised (``please bear with me'') instead of delivering. The same pattern held on the 81 dev-split cases (D 1.10 vs.\ C 0.57). Across all 241 evaluated cases spanning
all three DialMAGE splits (213 with a database request), REPAIR raises objective real, matching entity
delivery from 18.3\% (A), 20.6\% (B) and 17.4\% (C) to \textbf{62.9\%} (D) on the 8B model, and from 29.0\%
(A), 41.6\% (B) and 45.8\% (C) to \textbf{74.3\%} (D) on Claude Sonnet.

\paragraph{Strong model.} With Claude Sonnet as the generator (Table~\ref{tab:claude}), plain judges without
database verification favour prompt-only reading. However, fact-checked judging reverses the picture under
both judges: with Opus, D resolves 1.64 [1.52, 1.75] vs.\ C 0.82 and B 0.76 (D$-$C $+0.82$ [0.66, 0.97],
D$-$B $+0.88$ [0.71, 1.05]); with Gemini (cross-family), D resolves \textbf{1.94} [1.88, 1.99] vs.\ C 1.04
[0.90, 1.19] and B 0.91 [0.77, 1.05] (D$-$C $+0.90$ [0.76, 1.04], D$-$B $+1.03$ [0.89, 1.17]; D picked
best 53\% vs.\ C 6\%). REPAIR raises real, matching entities from 52.9\% (C) to 67.8\% and cuts entities that
violate the user's constraints from 33.3\% to 13.8\%. Plain judges reward confident invented bookings (see
below); fact-checked judges penalise them.

\begin{table}[h]
\centering
\begin{tabular}{lccccc}
\toprule
 & \multicolumn{2}{c}{Opus judge} & \multicolumn{2}{c}{Gemini judge} & \\
\cmidrule(lr){2-3} \cmidrule(lr){4-5}
Claude Sonnet & resolves & give up & resolves & give up & real \& match \\
\midrule
A last message & 0.83 & 35\% & 1.19 & 21\% & 31.0\% \\
B be concrete & 1.25 & 7\% & 1.62 & 2\% & 43.7\% \\
C prompt-only reading & 1.45 & 2\% & 1.85 & 0\% & 52.9\% \\
D REPAIR & 1.34 & 18\% & 1.55 & 14\% & 67.8\% \\
E ledger + database & 1.28 & 19\% & 1.49 & 14\% & 65.5\% \\
\midrule
A, fact-checked judge & 0.68 & 38\% & 0.89 & 24\% & \\
B, fact-checked judge & 0.76 & 19\% & 0.91 & 8\% & \\
C, fact-checked judge & 0.82 & 12\% & 1.04 & 0\% & \\
D, fact-checked judge & \textbf{1.64} & \textbf{3\%} & \textbf{1.94} & \textbf{1\%} & \\
E, fact-checked judge & 1.65 & 3\% & 1.90 & 1\% & \\
\bottomrule
\end{tabular}
\caption{Same 100 cases, Claude Sonnet generator. Plain judges reward invented bookings; fact-checked judges
penalise them. Gemini is a cross-family judge for both generator and first judge.}
\label{tab:claude}
\end{table}

\paragraph{Small model with REPAIR vs.\ strong model.} On the same 100 cases, Llama 8B with REPAIR is close
to Claude Sonnet told to be concrete under both judges (Opus: 1.22 vs.\ 1.25; Gemini: 1.51 vs.\ 1.62) and
above Claude answering the last message (Opus: 0.83; Gemini: 1.19).

\paragraph{What carries the gain.} D and E do not differ significantly on the 8B model overall (D$-$E
resolves: Opus $+0.06$ [$-0.07$, 0.19]; Gemini $+0.12$ [$-0.05$, 0.28]): on the 88 stage-2 cases the
explicit ledger and the database account for the entire gain. However, on the 12 stage-3 cases --- where the
user shows stronger repair escalation --- D significantly outperforms E on the 8B model: Opus D$-$E $+0.50$
[0.08, 0.92] (D 1.50 vs.\ E 1.00), Gemini $+0.67$ [0.25, 1.17] (D 1.92 vs.\ E 1.25), and give-ups drop from
33\% (E) to 0--8\% (D). For instance, when a user has asked three times for a 3-star hotel in the east (none
exists in the database, case DMAGE271\#8), E stalls with ``Let me check if there are any hotels,'' whereas D's
stage-3 policy immediately delivers the closest real match (Express by Holiday Inn, 2 stars in the east, with
address and phone), scoring 2 under both judges against 0 for E. The escalation-aware policy helps precisely
where users are most frustrated.

\paragraph{Judges reward invented bookings.} Across all 84 Sonnet replies that named a database entity
violating the user's constraints (\texttt{named\_invalid}), plain judges still awarded full marks
($\emph{resolves}=2$) 38\% of the time under Opus and 65\% under Gemini; fact-checked judging cuts this to
17\% and 21\% (and on the 8B run from 16\% and 44\% down to 10\% under both judges). For example, without the
database, both judges gave full marks to a reply that ``booked'' train TR0899 for a Leicester--Cambridge
request; in the database TR0899 runs Cambridge--London. Showing the judge the database rows removes this: for
the 8B model, B falls from 0.45 to 0.17 (Opus) and from 0.98 to 0.26 (Gemini), while D is 1.41 (Opus) and
1.49 (Gemini). Because the rows come from the same reader REPAIR uses, we report both passes.

\section{Limitations}
Evaluation uses two LLM judge families (Claude Opus and Gemini) plus an objective entity check; we do not
include a separate human rating study on the generated replies, relying instead on EmoWOZ's human emotion
annotations, DBDC3's 30-annotator breakdown labels, and exact database matching. The complaint-word list is
lexical: on 200 sampled negative turns, 7.0\% matched with 92.9\% precision (1 false positive from task
language ``no, I need a train...''); the low recall confirms our claim that word-level sentiment misses most
dissatisfaction. The repeat threshold (Jaccard $\ge 0.5$, yielding 46.8\%) is robust: 0.4 gives 53.7\%, 0.6
gives 37.0\%. The REPAIR reader identified the pending request in 19/20 spot-checked cases and provided
database rows in 13/20. Our corpus findings replicate across four datasets (EmoWOZ-MultiWOZ, EmoWOZ-DialMAGE,
USS-SGD across 20 domains, and USS-ReDial), with DBDC3 as a chat-oriented contrast; the database-grounded
intervention is evaluated on DialMAGE using the static MultiWOZ database and simulated booking. The objective
check depends on the reader's own reading of the pending request.

\section{Ethics}
EmoWOZ is released under CC BY-NC 4.0 \citep{feng2022emowoz}; we use it for academic research and do not
redistribute the data. All generated replies are from dialogue systems, not real users. No personal
information beyond what is in the public corpus was used.

\section{Conclusion}
Users tell a system it has failed mostly by repeating themselves, and systems answer the repetition instead of
the need behind it. Making that need explicit --- what was asked, what was delivered, what is still pending,
with real data to deliver it --- lets a small model do what a strong model does from a prompt, and makes both
deliver real, correct entities more often.

\paragraph{Reproducibility.} Code, prompts, model outputs and judgments:\\
\url{https://github.com/yingchen-coding/frustration-paper}.

\bibliographystyle{plainnat}
\bibliography{refs}
\end{document}
